\usepackage[provide=*]{babel}
\usepackage{datetime2}
\DTMsetup{useregional}

\usepackage[table,dvipsnames]{xcolor}
\definecolor{ecnblue}{RGB}{16,28,72}
\definecolor{ecnyellow}{RGB}{250,182,0}
\colorlet{titlecolor}{ecnblue}
\colorlet{secondarycolor}{ecnyellow}
\definecolor{insared}{RGB}{238,28,34}
\definecolor{insagrey}{RGB}{79,77,80}

% For tikz and pgf plots
\usepackage{tikz}
\usetikzlibrary{calc}
\usepackage{tikz-cd}
\usepackage{progressbar}
\usetikzlibrary{shapes.geometric}
\usetikzlibrary{babel}
\usepackage{pgfplots}
\usepackage{import}
\pgfplotsset{compat=1.14}
\def\mathdefault#1{#1} 			% avoid pgf bug from plt figure export
\usepackage{xsavebox} 			% save content for repeated use
\usepackage{atbegshi} 			% insert material on every page
\usepackage[texcoord=true]{eso-pic} % The zero point of the coordinate system is the upper left corner of the page.

% Maths set up 
\usepackage[retainorgcmds]{IEEEtrantools} % utlime equation envir. (see guide in LateX folder on Chrome)
\usepackage[intlimits,  	% places super/subscript on top of the integral sign
			reqno,      	% Place equation numbers on the right.
			]{amsmath} 
\usepackage{amsfonts}
\usepackage{amssymb}		% nice math symbols
%\usepackage{esvect} 		% vector sign tuning : \vv{e}{1} does not put arrow over subscript \vv{e_1} does
\usepackage{exscale} 		% to get the sum sign right (bigger)
\usepackage{xfrac} 			% nice inline fraction with \sfrac command

% Physics and chemistry 
\usepackage{siunitx}
\usepackage[version=4]{mhchem}
%\usepackage{chemfig}
%\usepackage{chemformula}

% Code
\usepackage{piton}           % shows code nicely
\usepackage{tcolorbox}       % box the pyhton code in a frey shaded box
\tcbuselibrary{breakable}    % idk

%\DeclareMathOperator{\logd}{$\mathrm{log}_2$} % example of the command \logd in math mode which write log_2 in upwrite nicely
\newcommand{\vect}[1]{\vec{#1}}  % example of a new command \vect for writing bold vectors
\newcommand{\overbar}[1]{\mkern1.5mu\overline{\mkern-1.5mu#1\mkern-1.5mu}\mkern 1.5mu}


% Bibliography 
\usepackage[sorting=none,      		% sorts by citation order
			citestyle=numeric-comp, % number style and concatenation
			autocite=superscript,
			language=french,
			style=nature,
			]{biblatex}
\usepackage{doi}

% Fonts

%lettrine stuff
\usepackage{lettrine,verbatim,longtable,url}
\usepackage{Konanur,Sanremo,Kramer,Typocaps}
\renewcommand{\LettrineFontHook}{\Typocapsfamily\color{insagrey}}
\renewcommand{\LettrineTextFont}{\color{insagrey}\scshape}

%first set up the main font
\usepackage{fontspec}
\setmainfont[Numbers = {OldStyle,Proportional},Ligatures=TeX,Microtype]{TeXGyrePagella}
\setsansfont[Microtype,Scale=MatchLowercase]{TeXGyreAdventor}
\newfontfamily\titlefamily{TeXGyreAdventor}[Microtype,Scale=MatchLowercase,LetterSpace=3]

\newfontfamily{\liningmainfont}[Numbers=Lining]{TeXGyrePagella}
\newfontface\heroscondensedbold{texgyreheroscn-bold.otf}[RawFeature = {+smcp}]
\newfontface\heroscondensedreg{texgyreheroscn-regular.otf}

%then set up the math font
\usepackage{unicode-math}
\setmathfont{TeX Gyre Pagella Math}

\linespread{1.05} 							% increase line spacing slightly
\usepackage[babel=true,nopatch=footnote]{microtype} 	% 'subliminal refinements towards typographical perfection'
\usepackage{romannum}

% Tables
\usepackage{tabularx} 		% automatically loaded with memoir, can fix the table width 
\usepackage{multirow} 		% ?
\usepackage{makecell} 		% retour à la ligne dans une cellule 
\usepackage{booktabs} 		% automatically loaded with memoir Required for better horizontal rules in tables

% Float captions
\usepackage{setspace}
\usepackage[font=small,
			textfont = {rm},
			format = plain,
			figureposition = below,
			tableposition = top,
			hypcap=true 			% for commands in caption (cite, math?)
			]{caption}
\DeclareCaptionFont{tc}{\color{titlecolor}}
\DeclareCaptionLabelSeparator{CapSep}{~\ding{71}~}
\captionsetup{labelfont={tc,sf,sc},labelsep=CapSep}

%Pages geometry
\usepackage[%showframe
]{geometry}
%\usepackage{pdflscape}
%\usepackage{rotating}
%\usepackage[export]{adjustbox}

\usepackage{csquotes}
\usepackage{textcomp} 				% for \textmu
\usepackage[language=french]{lipsum}
\usepackage{datatool-base}      	% for bib entries with ISBN
\usepackage{graphicx}
\usepackage{flafter}				% float appears after \ref. in text
\usepackage{pifont}					% for quick symbols, see https://lesfichesabebert.fr/TeX/latex/Listes.html
\usepackage[shortlabels]{enumitem}	% tuning listings environments
\usepackage{ragged2e}
\usepackage{todonotes}				% insert to do notes via \todo{}
\usepackage{titletoc}				% for \minitoc at chapter calls
\usepackage{cuted}					% for onecolumn domain w/o pageskips
\setlength{\stripsep}{0ex}
\usepackage{hyphenat}
\usepackage{etoolbox}
%\usepackage{lua-visual-debug}

\usepackage{hyperref}
\hypersetup{
	frenchlinks = true, 	% Use small caps instead of color for links.
	plainpages = false,
	colorlinks = true,
	linktocpage=true,		% numbers are only hyperlink in toc
	linkcolor = titlecolor,
	citecolor = titlecolor,
	urlcolor = titlecolor,
	pdfkeywords = \keywords,
	pdfsubject = \subject,
	pdfauthor = {\author},
	pdftitle = {\thetitle},
	pdfnewwindow = true, 	% new window for link to a place out of the doc
	pdfpagelayout = OneColumn,
	pdfstartview = FitBH,
	unicode = true, 		% permettre d’utiliser des caractères non-latins dans les marque-pages du lecteur de PDF
}

% Glossary
\usepackage[xindy={codepage=utf8}, 	% adding utf8 so that it works with Lua
	toc=true, 				% add glossaries to the toc.
	translate=babel, 		% If babel has been loaded, the translate=babel option is switched on
	%automake, 				% usable in unrestricted mode...
	%abbreviations,
	acronym, 				% add the category for acronyms
	symbols, 				% add the category for symbols
	%hyperfirst = false,	% idk
	%stylemods = mcols, 		% column style + correct style
	%style = mcolindex, 	% two column in textwidth
	%postdot, 				% puts a dot behind each entry
	nonumberlist, 			% no page refs in the glossaries
]{glossaries-extra}


%Usage:
%1.  add entry (e.g.\newabbreviation{svm}{SVM}{support vector machine}) in glossary.tex
%2. refers the glossary item in document by \gls{item}
%3. compile with cmd: >> makeglossary main 
%4. compile with cmd: >> lualatex main
%5. compile again: >> makeglossary main ,for news items refered in document

%\glsfmttext{label} using glossary entries in headings or labels
